% Complete analytic statement about ordinary products, not the full core.
\section{Geometrically weighted differences for all nonrectangular binary pairs}
For $\alpha\ge1$ and $h\ge0$, let $A_{\alpha,h}=F_{1^\alpha2^h}$.

\begin{lemma}[A binary growth bound after one two-term window]
If $f_n$ counts any nonempty finite downset of compositions with
parts in $\{1,2\}$, and $g_n=f_n+f_{n-1}$, then
\[
                 3g_n\le5g_{n-1}\qquad(n\ge2).
\]
Consequently, for every $\lambda\ge5/3$ and every binary arm $B$,
the coefficients of $\lambda^j[z^j](1+z)q_B$ form a unimodal
sequence, where $q_B$ is the reverse arm polynomial.
\end{lemma}
\begin{proof}
Deleting the first part gives $f_n\le f_{n-1}+f_{n-2}$.
The usual two extensions of a composition give
$f_{n-1}\le2f_{n-2}$ for $n\ge2$; the rank-one case is immediate.
Writing $x=f_{n-1}$ and $y=f_{n-2}$, we obtain
\[
 3g_n\le3(2x+y)\le5(x+y)=5g_{n-1}.
\]
Thus $g_n/\lambda^n$ is nonincreasing from rank one onward.
Its initial values are $g_0=1,g_1=2$, so it is unimodal whether
$\lambda$ is below or above two. Reversal gives the last assertion.
\end{proof}

\begin{lemma}[Two-term windows preserve the weighted difference property]
Let $p$ be nonnegative and unimodal. If
$\lambda^j d^+(p)_j$ is unimodal for some $\lambda>0$, then
$\lambda^j d^+((1+z)p)_j$ is unimodal as well.
\end{lemma}
\begin{proof}
Through the first mode $m$ of $p$, the new positive differences
agree with $(1+z)d^+(p)$. At rank $m+1$ the actual positive
difference is between zero and $d^+(p)_m$, and every later
positive difference vanishes. After weighting, the full window
is convolution by $(1,\lambda)$, a log-concave kernel. Such
convolution preserves unimodality. Terminal truncation and
decreasing the final remaining coefficient also preserve it.
\end{proof}

\begin{theorem}\label{thm:binary-geometric-differences}
Let $\alpha,\beta\ge1$ and $h,k\ge1$. Write
\[
 q(z)=z^{\alpha+\beta+2h+2k}
          A_{\alpha,h}(1/z)A_{\beta,k}(1/z).
\]
For every real $\lambda\ge5/3$, the sequence
$\lambda^j\max(q_j-q_{j-1},0)$ is unimodal.
The constant $5/3$ is sharp for a uniform upper-ray assertion
of this form.
\end{theorem}
\begin{proof}
Induct on $h+k$. Orient the weights as $a=\alpha+2h\ge b=\beta+2k$.
The first reverse mode $s$ of the ordinary product satisfies $s<a$:
use first-half growth if $a>b$, and the extra midpoint rank for
two nonrectangular factors if $a=b$.

If $h\ge2$, deleting the heavier arm's final two leaves both
arms nonrectangular and gives
\[
 q=(1+z)q_0+z^a q_B.
\]
Through $s<a$ this agrees with the two-term window. Its positive
differences are a prefix of those of that window, followed by
zeros. The induction hypothesis and the window lemma apply.

If $h=1$, do not delete the last two. Instead write
\[
 q=[\alpha+1]R+z^a q_B,\qquad R=(1+z)q_B,
 \qquad a=\alpha+2.
\]
For $j\le\alpha$, the first differences of $q$ equal $R_j$.
The only further possible positive difference before $s<a$
is at $j=\alpha+1$, where it equals $R_j-R_0=R_j-1$.
Thus the positive differences are a terminal truncation of $R$,
possibly with their final coefficient decreased. The preceding
binary growth bound proves their weighted unimodality. This
supplies every terminal case of the induction.

For sharpness, take $A=A_{\alpha,1}$ with $\alpha\ge6$ and
$B=A_{1,2}$. The positive reverse differences of $AB$ are exactly
\[
                         (1,4,6,5,3,2,1).
\]
For every $3/2<\lambda<5/3$, the three entries
$5\lambda^3,3\lambda^4,2\lambda^5$ have a strict valley.
Hence no smaller constant can replace $5/3$ in the uniform
assertion for all larger $\lambda$.
\end{proof}

\begin{corollary}[A coefficient criterion for every binary pair]
Let both binary skeletons be nonempty and let $M\ge3$. In the
balanced strict-core band $|a-b|\le M-4$, put
$R=\lfloor(a+b+M)/2\rfloor$. If the ordinary reverse product
$q$ is log-concave through rank $R$, the core and both relative
polynomials are unimodal.
\end{corollary}
\begin{proof}
At $x=5/3$ the sum of the first three inverse powers is
$3/5+9/25+27/125=147/125>1$. Therefore $\lambda_M>5/3$
for $M\ge3$. Apply the geometric signed-window theorem locally.
The reverse core is $h_Mq+z^{a+M}q_B+z^{b+M}q_A$; balance puts
the two corrections after $R$. Forward first-half growth joins
the resulting prefix to a nonincreasing remaining reverse tail.
The usual balanced-band equivalence gives the relative statements.
\end{proof}

The prefix log-concavity hypothesis in this corollary remains a
separate requirement. The theorem removes the false global
log-concavity requirement on positive differences, but does not
by itself prove every binary core or the general one-descent
conjecture.

\begin{proposition}[Adjacent initial capacities]
If $|\alpha-\beta|\le1$, the positive reverse differences of
$A_{\alpha,h}A_{\beta,k}$ are log-concave for every $h,k\ge0$.
The bound one is sharp for a statement uniform over $h,k$.
\end{proposition}
\begin{proof}
Induct by deleting a two from a heavier nonrectangular arm.
Such an arm always exists unless both are intervals: any
nonrectangular arm outweighs a rectangular one under the stated
capacity restriction. The first reverse mode is strictly below
the deleted arm's old weight. Apply the two-term-window
log-concave-difference lemma and terminal truncation. The base
is a product of two intervals.

For the sharpness assertion take $\alpha=\beta+2$, $h=0$, $k=1$,
with $\beta\ge1$. Both arm weights equal $a=\beta+2$. The first
reverse mode of $[a+1]q_{\beta,1}$ is $a$, and its positive
differences through that rank are exactly $q_{\beta,1}$.
Its last three coefficients are $(2,1,1)$, whose central
log-concavity determinant equals $-1$.
\end{proof}
